% Generated draft: requires mathematical and source review.
\documentclass[11pt]{article}
\usepackage[margin=25mm]{geometry}
\usepackage{fontspec}
\setmainfont{Latin Modern Roman}
\usepackage{amsmath,amssymb,mathtools,unicode-math}
\setmathfont{Latin Modern Math}
\usepackage{enumitem,needspace,xurl,hyperref}
\setlength{\parindent}{0pt}
\setlength{\parskip}{5pt}
\setlength{\emergencystretch}{3em}
\newcommand{\N}{\mathbb{N}}
\newcommand{\Z}{\mathbb{Z}}
\newcommand{\Q}{\mathbb{Q}}
\newcommand{\R}{\mathbb{R}}
\newcommand{\C}{\mathbb{C}}
\newcommand{\F}{\mathbb{F}}
\newcommand{\A}{\mathbb{A}}
\newcommand{\PP}{\mathbb{P}}
\newcommand{\E}{\mathbb{E}}
\newcommand{\eps}{\varepsilon}
\newcommand{\dd}{\,\mathrm{d}}
\begin{document}
\section*{UnsolvedMath ID: 700859}
\textbf{Source problem code:} catalogue entry
\subsection{Definitions and mathematical statement}
Let \(p\) be a prime and \(G\) a finite group. Interpret an irreducible \(p\)-modular character \(\varphi\) as the irreducible Brauer character afforded by a simple module \(V\) over an algebraically closed field of characteristic \(p\); its degree is \(\varphi(1)=\dim V\). For a positive integer \(n\), define its \(p\)-part by
\[
 n_p=p^{v_p(n)},\qquad v_p(n)=\max\{a\in\Z_{\geq0}:p^a\mid n\}.
\]
No hypothesis that \(p\mid |G|\) is stated in the source.

For a fixed prime \(p\), a finite-valued bound of the requested form is a function
\[
 f_p:\{p^a:a\in\Z_{\geq0}\}\longrightarrow\R_{\geq0}
\]
such that, for every finite group \(G\) and every irreducible \(p\)-modular character \(\varphi\) of \(G\),
\[
 \varphi(1)_p\leq f_p(|G|_p).
\]
The requested problem is to determine whether such bounds exist and, where they do, determine the best possible bound, separately for
\[
 \text{(a) }p=2,\qquad \text{(b) }p\ne2.
\]
Under the pointwise interpretation of ``best possible,'' the target for fixed \(p\) is the extended-real extremal function
\[
 B_p(p^a)=\sup\bigl\{\varphi(1)_p:\ G\text{ is finite},\ |G|_p=p^a,\ \varphi\text{ is an irreducible }p\text{-modular character of }G\bigr\}.
\]
A finite bound exists precisely when these suprema are finite for every \(a\geq0\).

The source writes \(f\), without explicitly stating whether part (b) permits dependence on the odd prime \(p\). If a single function is intended, require instead a finite-valued function on \(\{p^a:p\text{ is an odd prime},\ a\geq0\}\), with the displayed inequality holding simultaneously for all odd primes. This quantifier qualification remains unresolved.

The supplied record labels the overall problem \emph{partially solved} and reports part (a) as solved negatively: no estimate of the requested kind exists for \(p=2\). This is a supplied source assessment, not an independently verified assertion here. Part (b) is retained as the remaining question in that record.
\subsection{Short English statement}
For irreducible modular characters of finite groups, seek the sharpest bound on the \(p\)-part of the character degree using only the \(p\)-part of the group order. The supplied source reports that no such bound exists for \(p=2\), and retains the odd-prime case as unresolved. Whether the odd-prime bound must be uniform in \(p\) is not explicit.
\subsection{Sources}
\begin{itemize}
\item Supplied local UnsolvedMath registry record, permanent ID \(700859\), exact source code \texttt{catalogue entry}, titled ``Kourovka Notebook Problem 12.40.'' The record attributes the problem to A. S. Kondratiev, proposed in 1992, and labels it \texttt{partially\_solved}.
\item \emph{The Kourovka Notebook (Unsolved Problems in Group Theory)}, problem 12.40, originating in the 12th issue (1992), as represented in the supplied 21st-issue edition (2026), \texttt{arXiv:1401.0300v48}. The supplied collection description names E. I. Khukhro and V. D. Mazurov as editors.
\item The supplied discussion attributes the negative resolution of part (a) to M. Kinyon and J. Urban, \emph{Preprint}, 2026, describing it as an AI solution confirmed by A. S. Kondratiev: \url{https://kourovkanotebookorg.wordpress.com/wp-content/uploads/2026/09/kourovka-mathematics.pdf}. No preprint title is supplied.
\item The supplied status review is dated \texttt{2026-10-16}, identifies its evidence as \texttt{maintained\_source}, and cites the October 2026 21st issue for the reported resolution of part (a). This assessment is supplied dated curation, not independently verified; the referenced sources have not been visited for this draft.
\end{itemize}
\end{document}
